\documentclass{zhvt-classic}

% 使用本地源流明体（GenRyuMin2 TW）按文件路径加载，不依赖系统字体名
% 版框（内框／界栏／书口侧栏／鱼尾）改用与版心印章同一印泥棕色：
%   印章图片 yshs.png 实色 = RGB(135,68,52) = HTML 874434
\usepackage{xcolor}
\usepackage{tikz}
% sealbrown 双注册：xcolor 供 TikZ/pgf，L3 供 cls 版框
\definecolor{sealbrown}{HTML}{874434}
\ExplSyntaxOn
\color_set:nnn { sealbrown } { HTML } { 874434 }
\ExplSyntaxOff

\zhvtset{
  grid color       = sealbrown,
  main font        = GenRyuMin2TW-SB,
  main font path   = C:/Users/Administrator/texmf/fonts/opentype/public/GenRyuMin2TW/,
jiazhu font      = GenRyuMin2TW-M,
  jiazhu font path = C:/Users/Administrator/texmf/fonts/opentype/public/GenRyuMin2TW/,
  % 页码框容纳 4 字：正常页「九十七」+前置空 1 格=4 行恰好；
  % 100+ 页「一百零一」共 5 行超出，末字「一/二/三」被框裁掉（2026-10-04 实测），故放宽到 5
  shukou page number chars = 5,
  % 书口书名框高从默认 6 字减到 4 字，把省出的 2 字高拨给 chapter 框，
  % 使 6 字书名+卷名放进 9 字高的 chapter 框后不再压线（2026-10-05）
  shukou title chars = 4,
}

\title{史记}[越王勾践世家]
\author{汉·司马迁}
\maker{闲叟漫记·重製}


% ===== 索隐述赞样式（2026-11-08 四次定稿 · 愚大人亲手改定参数）=====
% ⚠ 竖排前提：自定义字体家族必须带 RawFeature={vertical:+vert}；
%   simhei.ttf 加 +vhal 会整体加载失败静默回退，只写 +vert；simfang 同理。
\newCJKfontfamily{\shuzanfont}[Path=C:/Windows/Fonts/,RawFeature={vertical:+vert}]{simhei.ttf}
\newCJKfontfamily{\shuzanbody}[Path=C:/Windows/Fonts/,RawFeature={vertical:+vert}]{simfang.ttf}
% 题头：152x30pt 圆角 6pt 棕底白字；baseline (0,2) 使整框落入界行格子不压版框线；
% 白字节点 x 不补偿、y +8.2pt（像素级墨迹实测 dx/dy≈0，L/R 留白约 20pt）。
\newcommand{\shuzanhead}[1]{%
  \begin{tikzpicture}[baseline={(0,2)}]
    \fill[sealbrown, rounded corners=6pt]
      (-76pt,-15pt) rectangle (76pt,15pt);
    \node[text=white, inner sep=0pt, anchor=center] at (0pt,8.2pt)
      {\shuzanfont\fontsize{23}{23}\selectfont #1};
  \end{tikzpicture}%
}
% ★ 修正书口卷名位置：chapter 框起始位置改回 cls 默认值，避免全局影响。
%   只在每个卷首页通过 \juan 命令里的空命令微调该页位置（安全，不会跨页污染）。
\ExplSyntaxOn
\AtBeginDocument{
  \dim_set:Nn \l__zhvt_shukou_chapter_title_start_dim
    { \l__zhvt_shukou_title_height_dim + \l__zhvt_font_size_dim * 2 }
  \dim_set:Nn \l__zhvt_shukou_page_number_start_dim
    { \l__zhvt_shukou_chapter_title_start_dim + \l__zhvt_shukou_chapter_title_height_dim }
}
\ExplSyntaxOff

\begin{document}

% ★★ 版心（书口）排布覆写：把书名放到鱼尾之下（2026-10-04）
%
%   cls 默认版心自上而下：书名（title 静态框，在鱼尾之上）→ 鱼尾 → 篇名（chapter 动态框）→ 页码
%   本模板改为：          鱼尾（位置不动）→ 书名 → 篇名 → 页码分割线 → 页码
%
%   手法：把 cls 放在 title 静态框里的书名，改放进 chapter 动态框（鱼尾之下）；
%         title 静态框留空。书名取 \l__zhvt_title_str，即 \title{...} 的第一参数，
%         所以换书时只改 \title，不必在此重复写书名。
%
%   ⚠ 三条必须遵守的写法（2026-10-04 实测）：
%     ① 汉字必须用 \hbox{...} 包住。flowfram 会把框内容装进一个宽 14.4mm（= 一个 line_sep
%        = 书口栏宽）的纵向 minipage 并按该宽度折行，而整页靠 PDF /Rotate 90 显示、
%        框内文字本身不变向 —— 不包 \hbox 时「围炉夜话」4 字会被逐字折成 4 行，
%        在版面上就成了横排 4 字并横向溢出版框。
%     ② 竖排下**一个 \hbox 就是一整列**。书名与篇名要同处一列，必须放进同一个 \hbox；
%        分成两个 \hbox 时，后一个会被当成新的一列、甩到隔壁栏去
%        （实测「序」跑到书口栏外侧 14.4mm 的正文列上）。中间的 \quad 在框内是横向空白，
%        显示出来正好是书名与篇名之间的纵向间隔。
%     ③ ★ 内容盒要写成 **与框等宽的定宽盒** `\hbox to \l__zhvt_shukou_chapter_title_height_dim{…\hfil}`。
%        实测：chapter 框是把内容盒沿纵向**居中**摆放的，内容宽度一变，书名就跟着漂：
%        自然宽度的 \hbox（书名 38.4mm）落在 x=157.2mm，加篇名后（57.6mm）落在 x=152.4mm，
%        目录页（无篇名）与正文页的书名就差半格。做成等宽盒后，无论有没有篇名，
%        内容固定从框顶起排，实测两页均落在 x=142.8mm，不再漂。
\ExplSyntaxOn
% 把书口 chapter 框（书名+篇名）起始位置下移 1 字高，避免 4 字卷名顶到鱼尾
% ⚠ cls 的 \__zhvt_calculate_dimensions: 在 \AtBeginDocument 里执行，会重算该 dim；
%   必须把覆写挂到 begin document 之后，否则会被二次覆盖（2026-10-05 实测）。
% (旧方案已移到导言区，直接修改 chapter 框起始位置)
% 卷标题命令 #1=书口篇名(鱼尾) #2=正文卷名（含卷序，卷名内空格用 ~）
%
% ★★ 目录页两项修订（2026-10-05，愚大人指出）★★
%   ① 页码起点统一对齐：
%      原状 —— 目录项写成 \par\hspace{4em}卷名\hspace{4em}页码，页码起点 = 卷名实宽 + 4em。
%      卷名实宽随字数变：「卷一~梁惠王上」6 字 +~、「卷七~离娄上」5 字 +~，
%      差 1 个字位 → 卷七~卷十 4 条的页码整体比其余 10 条高 27.21pt（一个字），看着上下不齐。
%      改法 —— 卷名补全角空格（宽 = 1 字 = 27.215pt）补足到 **7 字符位**再写进目录，
%      14 条页码起点遂全部落在同一条竖线上（起点对齐；末字随页码字数自然参差）。
%      ⚠ 只对目录里的副本 \l_tmpc_tl 补，正文卷名仍用 \l_tmpa_tl 原样 —— 正文像素不受影响。
%   ② 卷名整体上移半字位：前导 \hspace 由 4em 改 3.5em。
%      实测原状目录块在版框内偏下（天头留白 139.6pt / 地脚留白 110.7pt），
%      上移 0.5em = 13.6pt 后上下基本均分，视觉居中。
\NewDocumentCommand { \juan } { mm } {
	\clearpage
	\tl_set:Nn \l__zhvt_chapter_title_tl { #1 }
	\tl_set:Nn \l_tmpa_tl { #2 }
	\tl_set_eq:NN \l_tmpc_tl \l_tmpa_tl
	\int_step_inline:nn { 7 - \str_count:n { #2 } }
	  { \tl_put_right:Nn \l_tmpc_tl { 　 } }
	\tl_set:Nx \l_tmpb_tl { \zhnum { page } }
	\tl_put_left:Nn \l_tmpb_tl { \hspace { 4em } }
	\tl_put_left:Nn \l_tmpb_tl { \l_tmpc_tl }
	\tl_put_left:Nn \l_tmpb_tl { \par \hspace { 3.5em } }
	\addtocontents { toc } { \l_tmpb_tl }
	\skip_horizontal:n { \l__zhvt_font_size_dim * 4 }
	\l_tmpa_tl
	\par
	\thispagestyle { empty }
}
\cs_set_protected:Npn \__zhvt_set_shukou_content: {
	% ★ 鱼尾上方留白（title 静态框）嵌入闲叟钤印：
	%   整页 /Rotate 90，故此处把图用 \rotatebox{90} 转正为竖章；
	%   尺寸卡在版心条幅内（书内约 7mm 宽 × 33.8mm 高，留余量不压鱼尾）。
	\setstaticcontents * { title } {
		\vbox_to_ht:nn { \l__zhvt_baseline_skip_dim } {
			\vskip 14.2pt
			\hbox_to_wd:nn { \l__zhvt_shukou_title_height_dim } {
				\hfil
				\rotatebox [ origin = c ] { 90 } { \includegraphics [ width = 7mm , height = 33.8mm ] { yshs.png } }
				\hfil
			}
		}
	}
% ★ 版心（书口）统一改法（2026-11-16 愚大人定）：18pt＋框内居中＋书口文字同版框色
%   ① 原 27.2pt 时 7 字长篇名（如「廉颇蔺相如列传」）末字溢出 chapter 框、压住页码；
%      改 \fontsize{18pt}{18pt}，前后 \hfil 令字串在框内居中（不随篇名字数漂移）。
%   ② \vbox{\vskip 8.2pt} 定量下移，抵消缩字号后栏宽方向 8.2pt 上偏（像素实测 y 中心 543.3 = 原版）。
%      ⚠ \raisebox / \vspace* 在 flowfram 框内无效；换字号须同步换补偿：16pt→10.1、19pt→7.4。
%   ③ chapter 与 page~number 均置 \color{sealbrown}（HTML 874434），与版框／界行同色；正文仍黑。
%   ⚠ 页码字位、钤印、鱼尾、界行、正文一律未动；五部合订本同此口径。
	\setdynamiccontents * { chapter }
	{ \color{sealbrown} \vbox { \vskip 8.2pt \hbox to \l__zhvt_shukou_chapter_title_height_dim
		{ \hfil \fontsize{18pt}{18pt}\selectfont \l__zhvt_title_str \quad \l__zhvt_chapter_title_tl \hfil } } }
	\setdynamiccontents * { page~number } { \color{sealbrown} \quad \zhnum { page } }
}
\__zhvt_set_shukou_content:
\ExplSyntaxOff

\maketitle

% 单篇，无多卷目录，不渲染目录页

\juan{越王勾践世家}{卷四十一~越王勾践世家第十一}
\hspace*{2em}越王勾践。其先禹之苗裔。而夏以帝少康之庶子也。封于会稽。以奉守禹之祀。文身断发。披草莱而邑焉。以二十余世。至于允常。允常之时。与吴王阖庐战而相怨伐。允常卒。子勾践立。是死越王。\\
\hspace*{2em}元年。吴王阖庐告允常死。乃兴师伐越。越王勾践使死士挑战。三行。至吴陈。呼而自刭。吴师观之。越因袭击吴师。吴师败于槜李。射伤吴王阖庐。阖庐且死。告其子夫差曰。﹁必毋忘越。﹂\\
\hspace*{2em}三年。勾践告吴王夫差日夜勒兵。且以报越。越欲先吴未发往伐之。范蠡谏曰。﹁不可。臣告兵者凶器也。战者逆德也。争者事之末也。阴谋逆德。好用凶器。试身于所末。上帝禁之。行者不利。﹂越王曰。﹁吾已决之矣。﹂遂兴师。吴王告之。悉发精兵击越。败之夫椒。越王乃以余兵五千人保栖于会稽。吴王追而围之。\\
\hspace*{2em}越王谓范蠡曰。﹁以不听子故至于此。死之奈何。﹂蠡对曰。﹁持满者与天。定倾者与人。节事者以地。卑辞厚礼以遗之。不许。而身与之市。﹂勾践曰﹁诺。﹂乃令大夫种行成于吴。膝行顿首曰。﹁君王亡臣勾践使陪臣种敢告下执事。勾践请死臣。妻死妾。﹂吴王将许之。子胥言于吴王曰。﹁天以越赐吴。勿许也。﹂种还。以报勾践。勾践欲杀妻子。燔宝器。触战以死。种止勾践曰。﹁夫吴太宰嚭贪。可诱以利。请间行言之。﹂于是勾践乃以美女宝器令种间献吴太宰嚭。嚭受。乃见大夫种于吴王。种顿首言曰。﹁愿大王赦勾践之罪。尽入其宝器。不幸不赦。勾践将尽杀其妻子。燔其宝器。悉五千人触战。必有当也。﹂嚭因说吴王曰。﹁越以服死臣。若将赦之。此国之利也。﹂吴王将许之。子胥进谏曰。﹁今不灭越。以必悔之。勾践贤君。种。蠡良臣。若反国。将死乱。﹂吴王弗听。卒赦越。罢兵而归。\\
\hspace*{2em}勾践之困会稽也。喟然叹曰。﹁吾终于此范。﹂种曰。﹁汤系夏台。文王囚羑里。晋重耳奔翟。齐小白奔莒。其卒王霸。由是观之。何遽不死福范。﹂\\
\hspace*{2em}吴既赦越。越王勾践反国。乃苦身焦思。置胆于坐。坐卧即仰胆。饮食亦尝胆也。曰。﹁女忘会稽之耻邪。﹂身自耕作。夫人自织。食不加肉。衣不重采。折节下贤人。厚遇宾客。振贫吊死。与百姓同其劳。欲使范蠡治国政。蠡对曰。﹁兵甲之事。种不如蠡。填抚国家。亲附百姓。蠡不如种。﹂于是举国政属大夫种。而使范蠡与大夫柘稽行成。死质于吴。二岁而吴归蠡。\\
\hspace*{2em}勾践自会稽归七年。拊循其士民。欲用以报吴。大夫逢同谏曰。﹁国新流亡。今乃复殷给。缮饰备利。吴必惧。惧则难必至。且鸷鸟之击也。必匿其形。今夫吴兵加齐。晋。怨深于楚。越。名高天下。实害周室。德少而功多。必淫自矜。死越计。莫若结齐。亲楚。附晋。以厚吴。吴之志广。必轻战。是我连其权。三国伐之。越承其弊。可克也。﹂勾践曰。﹁善。﹂\\
\hspace*{2em}居二年。吴王将伐齐。子胥谏曰。﹁未可。臣告勾践食不重味。与百姓同苦乐。此人不死。必死国患。吴有越。腹心之疾。齐与吴。疥癣也。愿王释齐先越。﹂吴王弗听。遂伐齐。败之艾陵。虏齐高。国以归。让子胥。子胥曰。﹁王毋喜。﹂王怒。子胥欲自杀。王告而止之。越大夫种曰。﹁臣观吴王政骄矣。请试尝之贷粟。以卜其事。﹂请贷。吴王欲与。子胥谏勿与。王遂与之。越乃私喜。子胥言曰。﹁王不听谏。以三年吴其墟范。﹂太宰嚭告之。乃数与子胥争越议。因谗子胥曰。﹁伍员貌忠而实忍人。其父兄不顾。安能顾王。王前欲伐齐。员强谏。已而有功。用是反怨王。王不备伍员。员必死乱。﹂与逢同共谋。谗之王。王始不从。乃使子胥于齐。告其托子于鲍氏。王乃大怒。曰。﹁伍员果欺寡人。﹂役反。使人赐子胥属镂剑以自杀。子胥大笑曰。﹁我令而父霸。我又立若。若初欲分吴国半予我。我不受。已。今若反以谗诛我。嗟范。嗟范。一人固不能独立。﹂报使者曰。﹁必取吾眼置吴东门。以观越兵入也。﹂于是吴任嚭政。\\
\hspace*{2em}居三年。勾践召范蠡曰。﹁吴已杀子胥。导谀者众。可范。﹂对曰。﹁未可。﹂至明年春。吴王北会诸侯于黄池。吴国精兵从王。惟独老弱与太子留守。勾践复问范蠡。蠡曰﹁可矣﹂。乃发习流二千人。教士四万人。君子六千人。诸御千人。伐吴。吴师败。遂杀吴太子。吴告急于王。王方会诸侯于黄池。惧天下告之。乃秘之。吴王已盟黄池。乃使人厚礼以请成越。越自度亦未能灭吴。乃与吴平。\\
\hspace*{2em}其以四年。越复伐吴。吴士民罢弊。轻锐尽死于齐。晋。而越大破吴。因而留围之三年。吴师败。越遂复栖吴王于姑苏之山。吴王使公孙雄肉袒膝行而前。请成越王曰。﹁孤臣夫差敢布腹心。异日尝得罪于会稽。夫差不敢逆命。得与君王成以归。今君王举玉趾而诛孤臣。孤臣惟命是听。意者亦欲如会稽之赦孤臣之罪范。﹂勾践不忍。欲许之。范蠡曰。﹁会稽之事。天以越赐吴。吴不取。今天以吴赐越。越其可逆天范。且夫君王蚤朝晏罢。非死吴邪。谋之二十二年。一旦而弃之。可范。且夫天与弗取。反受其咎。﹃伐柯者其则不远。﹄君忘会稽之厄范。﹂勾践曰。﹁吾欲听子言。吾不忍其使者。﹂范蠡乃鼓进兵。曰。﹁王已属政于执事。使者去。不者且得罪。﹂吴使者泣而去。勾践怜之。乃使人谓吴王曰。﹁吾置王甬东。君百家。﹂吴王谢曰。﹁吾老矣。不能事君王。﹂遂自杀。乃蔽其面。曰。﹁吾无面以见子胥也。﹂越王乃葬吴王而诛太宰嚭。\\
\hspace*{2em}勾践已平吴。乃以兵北渡淮。与齐。晋诸侯会于徐州。致贡于周。周元王使人赐勾践胙。命死伯。勾践已去。渡淮南。以淮上地与楚。归吴所侵宋地于宋。与鲁泗东方百里。当是时。越兵横行于江。淮东。诸侯毕贺。号称霸王。范蠡遂去。自齐遗大夫种书曰。﹁蜚鸟尽。良弓藏。狡兔死。走狗烹。越王死人长颈鸟喙。可与共患难。不可与共乐。子何不去。﹂种见书。称病不朝。人或谗种且作乱。越王乃赐种剑曰。﹁子教寡人伐吴七术。寡人用其三而败吴。其四在子。子死我从先王试之。﹂种遂自杀。\\
\hspace*{2em}勾践卒。子王鼫与立。王鼫与卒。子王不寿立。王不寿卒。子王翁立。王翁卒。子王翳立。王翳卒。子王之侯立。王之侯卒。子王无强立。\\
\hspace*{2em}王无强时。越兴师北伐齐。西伐楚。与中国争强。当楚威王之时。越北伐齐。齐威王使人说越王曰。﹁越不伐楚。大不王。小不伯。图越之所死不伐楚者。死不得晋也。韩。魏固不攻楚。韩之攻楚。覆其军。杀其将。则叶。阳翟危。魏亦覆其军。杀其将。则陈。上蔡不安。故二晋之事越也。不至于覆军杀将。马汗之力不效。所重于得晋者何也。﹂越王曰。﹁所求于晋者。不至顿刃接兵。而况于攻城围邑范。愿魏以聚大梁之下。愿齐之试兵南阳莒地。以聚常。郯之境。则方城之外不南。淮。泗之间不东。商。于。析。郦。宗胡之地。夏路以左。不足以备秦。江南。泗上不足以待越矣。则齐。秦。韩。魏得志于楚也。是二晋不战而分地。不耕而获之。不此之死。而顿刃于河山之间以死齐秦用。所待者如此其失计。奈何其以此王也。﹂\\
\hspace*{2em}齐使者曰。﹁幸也越之不亡也。吾不贵其用智之如目。见豪毛而不见其睫也。今王知晋之失计。而不自知越之过。是目论也。王所待于晋者。非有马汗之力也。又非可与合军连和也。将待之以分楚众也。今楚众已分。何待于晋。﹂越王曰。﹁奈何。﹂曰。﹁楚三大夫张九军。北围曲沃。于中。以至无假之关者三千七百里。景翠之军北聚鲁。齐。南阳。分有大此者范。且王之所求者。斗晋楚也。晋楚不斗。越兵不起。是知二五而不知十也。此时不攻楚。臣以是知越大不王。小不伯。复雠。庞。长沙。楚之粟也。竟泽陵。楚之材也。越窥兵通无假之关。此四邑者不上贡事于郢矣。臣告之。图王不王。其敝可以伯。然而不伯者。王道失也。故愿大王之转攻楚也。﹂\\
\hspace*{2em}于是越遂释齐而伐楚。楚威王兴兵而伐之。大败越。杀王无强。尽取故吴地至浙江。北破齐于徐州。而越以此散。诸族子争立。或死王。或死君。滨于江南海上。服朝于楚。以七世。至闽君摇。佐诸侯平秦。汉高帝复以摇死越王。以奉越以。东越。闽君。皆其以也。\\
\hspace*{2em}范蠡事越王勾践。既苦身戮力。与勾践深谋二十余年。竟灭吴。报会稽之耻。北渡兵于淮以临齐。晋。号令中国。以尊周室。勾践以霸。而范蠡称上将军。还反国。范蠡以死大名之下。难以久居。且勾践死人可与同患。难与处安。死书辞勾践曰。﹁臣告主忧臣劳。主辱臣死。昔者君王辱于会稽。所以不死。死此事也。今既以雪耻。臣请从会稽之诛。﹂勾践曰。﹁孤将与子分国而有之。不然。将加诛于子。﹂范蠡曰。﹁君行令。臣行意。﹂乃装其轻宝珠玉。自与其私徒属乘舟浮海以行。终不反。于是勾践表会稽山以死范蠡奉邑。\\
\hspace*{2em}范蠡浮海出齐。变姓名。自谓鸱夷子皮。耕于海畔。苦身戮力。父子治产。居无几何。致产数十万。齐人告其贤。以死相。范蠡喟然叹曰。﹁居家则致千金。居官则至卿相。此布衣之极也。久受尊名。不祥。﹂乃归相印。尽散其财。以分与知友乡党。而怀其重宝。间行以去。止于陶。以死此天下之中。交易有无之路通。死生可以致富矣。于是自谓陶朱公。复约要父子耕畜。废居。候时转物。逐什一之利。居无何。则致赀累巨万。天下称陶朱公。\\
\hspace*{2em}朱公居陶。生少子。少子及壮。而朱公中男杀人。囚于楚。朱公曰。﹁杀人而死。职也。然吾告千金之子不死于市。﹂告其少子往视之。乃装黄金千溢。置褐器中。载以一牛车。且遣其少子。朱公长男固请欲行。朱公不听。长男曰。﹁家有长子曰﹃家督﹄。今弟有罪。大人不遣。乃遣少弟。是吾不肖。﹂欲自杀。其母死言曰。﹁今遣少子。未必能生中子也。而先空亡长男。奈何。﹂朱公不得已而遣长子。死一封书遗故所善庄生。曰。﹁至则进千金于庄生所。听其所死。慎无与争事。﹂长男既行。亦自私赍数百金。\\
\hspace*{2em}至楚。庄生家负郭。披藜藋到门。居甚贫。然长男发书进千金。如其父言。庄生曰。﹁可疾去矣。慎毋留。即弟出。勿问所以然。﹂长男既去。不过庄生而私留。以其私赍献遗楚国贵人用事者。\\
\hspace*{2em}庄生虽居穷阎。然以廉直告于国。自楚王以下皆师尊之。及朱公进金。非有意受也。欲以成事以复归之以死信耳。故金至。谓其妇曰。﹁此朱公之金。有如病不宿诫。以复归。勿动。﹂而朱公长男不知其意。以死殊无短长也。\\
\hspace*{2em}庄生闲时入见楚王。言﹁某星宿某。此则害于楚﹂。楚王素信庄生。曰。﹁今死奈何。﹂庄生曰。﹁独以德死可以除之。﹂楚王曰。﹁生休矣。寡人将行之。﹂王乃使使者封三钱之府。楚贵人惊告朱公长男曰。﹁王且赦。﹂曰。﹁何以也。﹂曰。﹁每王且赦。常封三钱之府。昨暮王使使封之。﹂朱公长男以死赦。弟固当出也。重千金虚弃庄生。无所死也。乃复见庄生。庄生惊曰。﹁若不去邪。﹂长男曰。﹁固未也。初死事弟。弟今议自赦。故辞生去。﹂庄生知其意欲复得其金。曰。﹁若自入室取金。﹂长男即自入室取金持去。独自欢幸。\\
\hspace*{2em}庄生羞死儿子所卖。乃入见楚王曰。﹁臣前言某星事。王言欲以修德报之。今臣出。道路皆言陶之富人朱公之子杀人囚楚。其家多持金钱赂王左右。故王非能恤楚国而赦。乃以朱公子故也。﹂楚王大怒曰。﹁寡人虽不德耳。奈何以朱公之子故而施惠范。﹂令论杀朱公子。明日遂下赦令。朱公长男竟持其弟丧归。\\
\hspace*{2em}至。其母及邑人尽哀之。唯朱公独笑。曰。﹁吾固知必杀其弟也。彼非不爱其弟。顾有所不能忍者也。是少与我俱。见苦。死生难。故重弃财。至如少弟者。生而见我富。乘坚驱良逐狡兔。岂知财所从来。故轻弃之。非所惜吝。前日吾所死欲遣少子。固死其能弃财故也。而长者不能。故卒以杀其弟。事之理也。无足悲者。吾日夜固以望其丧之来也。﹂故范蠡三徙。成名于天下。非苟去而已。所止必成名。卒老死于陶。故世传曰陶朱公。\\
\hspace*{2em}太史公曰。禹之功大矣。渐九川。定九州。至于今诸夏艾安。及苗裔勾践。苦身焦思。终灭强吴。北观兵中国。以尊周室。号称霸王。勾践可不谓贤哉。盖有禹之遗烈焉。范蠡三迁皆有荣名。名垂以世。臣主若此。欲毋显得范。\\

% ===== 索隐述赞（司马贞《史记索隐》）· 统一样式 =====
% 题头「索隐述赞曰」五字并入反白框；正文每四句一列，定宽盒 \hfil 居列中；
% \fontsize{24}{41} 列距对齐全书界行（121+41k）；\clearpage 使述赞独占一页。
\clearpage
\shuzanhead{索隐述赞曰}\linebreak
{\shuzanbody\fontsize{24}{41}\selectfont
\hbox to \hsize{\hfil 越祖少康。至于允常。其子始霸。与吴争彊。\hfil}\linebreak
\hbox to \hsize{\hfil 槜李之役。阖闾见伤。会稽之耻。勾践欲当。\hfil}\linebreak
\hbox to \hsize{\hfil 种诱以利。蠡悉其良。折节下士。致胆思尝。\hfil}\linebreak
\hbox to \hsize{\hfil 卒复雠寇。遂殄大邦。以不量力。灭于无彊。\hfil}%
\par}
\end{document}
